\documentclass[11pt,letterpaper]{article}
\usepackage[margin=0.65in]{geometry}
\usepackage[T1]{fontenc}
\usepackage{lmodern,microtype,booktabs,tabularx,enumitem,xcolor,fancyhdr}
\usepackage[colorlinks=true,urlcolor=NDBlue,linkcolor=NDBlue]{hyperref}
\definecolor{NDBlue}{HTML}{0C2340}
\definecolor{NDGold}{HTML}{C99700}
\renewcommand{\familydefault}{\sfdefault}
\setlength{\parindent}{0pt}
\setlength{\parskip}{5pt}
\setlist[itemize]{leftmargin=1.2em,itemsep=2pt,topsep=3pt}
\setlist[enumerate]{leftmargin=1.5em,itemsep=3pt,topsep=3pt}
\pagestyle{fancy}
\fancyhf{}
\fancyfoot[L]{\footnotesize Notre Dame CBE \quad 28 September / 12 October 2026}
\fancyfoot[R]{\footnotesize\thepage\ / 2}
\renewcommand{\headrulewidth}{0pt}
\newcommand{\topic}[1]{\par\vspace{5pt}{\large\bfseries\color{NDBlue}#1}\par}
\newcommand{\repo}{https://github.com/dowlinglab/claude-for-researchers}
\begin{document}
{\LARGE\bfseries\color{NDBlue}Claude for Research}\par
{\large Beyond the Chatbot: a working reference}\par
{\color{NDGold}\rule{\linewidth}{1pt}}

\textbf{Keep the project's memory in files.} The conversation helps you work;
your repository records the evidence, decisions, and next steps that let you
resume. An agent's confident explanation is a proposal to inspect.

\topic{A small cycle you can reuse}
\begin{enumerate}
  \item \textbf{Inspect before editing.} Give the agent the relevant files and a
  bounded question. Ask for its understanding and a plan before it changes them.
  \item \textbf{Save what currently works.} Run the existing computation. Capture
  its outputs, parameters, units, solver settings, environment, and producing
  command. Commit the baseline before changing the model.
  \item \textbf{Change one thing.} Use a branch. Review the proposed scope and the
  resulting diff. Separate restructuring from new scientific assumptions.
  \item \textbf{Check against independent evidence.} Compare with the saved
  output or a source you opened. A test generated from the changed code can
  repeat the same mistake; explain where its expected value came from.
  \item \textbf{Leave a handoff.} Record what ran, what failed or remains unchecked,
  the evidence commit, and the next concrete action. Resume from this file.
\end{enumerate}

\topic{What belongs in the project}
\renewcommand{\arraystretch}{1.12}
\begin{tabularx}{\linewidth}{@{}p{1.72in}X@{}}
\toprule
\textbf{Artifact} & \textbf{What it should let someone do}\\
\midrule
Instructions file & Find the project-specific units, conventions, commands, and boundaries. Keep it short; point to detail elsewhere.\\
Environment + README & Build the environment and run the documented command from the stated directory.\\
Baseline + manifest & Trace a result to its inputs, settings, software versions, and original output. Preserve it when new output differs.\\
Research log & Recover decisions, rejected approaches, and the evidence that changed your mind.\\
Handoff & See checked and unchecked work, the branch and evidence commit, and a runnable next step.\\
\bottomrule
\end{tabularx}

\topic{The two-session project}
\textbf{September 28:} write project instructions, run the untouched notebook,
and capture a baseline. Finish by saving the state you actually reached.

\textbf{October 12:} resume extraction with one function, then audit selected
claims in the shared report. Switch to the audit even if migration is incomplete;
the unchanged notebook supplies the evidence.

Each meeting: \textbf{35 min presentation / 10 Q\&A / 45 hands-on / 15 regroup}.
Setup is pre-work. Between meetings, experiment on your own research.
Extraction resumes in Session 2. Follow-up can finish a baseline check or a small
claim set; the longer migration and writing workflows are extensions.

\newpage
{\Large\bfseries\color{NDBlue}A claim needs evidence you can locate}\par
{\color{NDGold}\rule{\linewidth}{1pt}}

\topic{A practical audit}
\begin{enumerate}
  \item Quote the claim and record its location before judging it.
  \item State what evidence would be required. Then locate what actually exists.
  \item Open the artifact or source yourself. For numbers, inspect the producing
  output. For literature, check the source's conditions and what it actually says.
  \item Choose a verdict and write the supported wording. Preserve the original
  quotation so the revision remains reviewable.
\end{enumerate}

\textbf{Illustrative example, unrelated to the CSTR exercise:}\par
\begin{tabularx}{\linewidth}{@{}p{1.28in}X@{}}
\toprule
Claim & ``All figures regenerate from a clean checkout.''\\
Required evidence & A clean-checkout run of the commands for every figure.\\
Actual evidence & A saved run log covers Figure A; the other figures have no run record.\\
Verdict & Cannot verify the full claim. Figure A has evidence.\\
Revision & ``Figure A regenerated in the recorded environment. The remaining figures have not yet been checked.''\\
\bottomrule
\end{tabularx}

\textbf{Useful verdicts:} verified; supported with limitations; mismatch; cannot
verify; not comparable. Missing evidence is a finding. A source about a different
system may be relevant without supporting a numerical comparison. A sample of
checked claims does not establish that the entire report is correct.

\topic{A prompt you can adapt}
\begin{quote}\small
Read these files before editing: [paths]. My question is [question]. List the
claims or changes you propose and the evidence each needs. Separate what you
observed from what you inferred. Show exact source locations. Stop at [boundary]
and save unresolved questions in [file]. I will review the evidence and diff.
\end{quote}

\topic{Before you stop working}
\begin{itemize}
  \item Inspect the diff and working tree. Commit only the changes you reviewed.
  \item Record commands and outcomes, including failures and checks you did not run.
  \item Keep the baseline unchanged when a comparison fails. Investigate the difference.
  \item Write the next command or question in the handoff. Identify the evidence commit
  it describes; commit the handoff separately.
  \item Check current institutional and sponsor rules before sharing research data
  with a tool. Use the seminar's references to find the authoritative guidance.
\end{itemize}

\topic{Find the reusable resources}
\small\raggedright
\href{\repo}{\textbf{github.com/dowlinglab/claude-for-researchers}}\\
In the repository: \path{resources/workshops/cstr/activities/session_guide.md}
for the in-room route; \texttt{resources/templates/} for project instructions,
a research log, and result provenance; \texttt{resources/practices/} for the
computing, literature, and manuscript workflows.\par
The claim-audit template is in
\path{resources/workshops/cstr/docs/claim_evidence_audit.md}.
Current policy and product sources are indexed in \texttt{notes/references.md}.
\end{document}
